%% Generated by Sphinx.
\def\sphinxdocclass{report}
\documentclass[letterpaper,10pt,italian]{sphinxmanual}
\ifdefined\pdfpxdimen
   \let\sphinxpxdimen\pdfpxdimen\else\newdimen\sphinxpxdimen
\fi \sphinxpxdimen=.75bp\relax
\ifdefined\pdfimageresolution
    \pdfimageresolution= \numexpr \dimexpr1in\relax/\sphinxpxdimen\relax
\fi
%% let collapsible pdf bookmarks panel have high depth per default
\PassOptionsToPackage{bookmarksdepth=5}{hyperref}

\PassOptionsToPackage{booktabs}{sphinx}
\PassOptionsToPackage{colorrows}{sphinx}

\PassOptionsToPackage{warn}{textcomp}
\usepackage[utf8]{inputenc}
\ifdefined\DeclareUnicodeCharacter
% support both utf8 and utf8x syntaxes
  \ifdefined\DeclareUnicodeCharacterAsOptional
    \def\sphinxDUC#1{\DeclareUnicodeCharacter{"#1}}
  \else
    \let\sphinxDUC\DeclareUnicodeCharacter
  \fi
  \sphinxDUC{00A0}{\nobreakspace}
  \sphinxDUC{2500}{\sphinxunichar{2500}}
  \sphinxDUC{2502}{\sphinxunichar{2502}}
  \sphinxDUC{2514}{\sphinxunichar{2514}}
  \sphinxDUC{251C}{\sphinxunichar{251C}}
  \sphinxDUC{2572}{\textbackslash}
\fi
\usepackage{cmap}
\usepackage[T1]{fontenc}
\usepackage{amsmath,amssymb,amstext}
\usepackage{babel}



\usepackage{tgtermes}
\usepackage{tgheros}
\renewcommand{\ttdefault}{txtt}



\usepackage[Sonny]{fncychap}
\ChNameVar{\Large\normalfont\sffamily}
\ChTitleVar{\Large\normalfont\sffamily}
\usepackage{sphinx}

\fvset{fontsize=auto}
\usepackage{geometry}


% Include hyperref last.
\usepackage{hyperref}
% Fix anchor placement for figures with captions.
\usepackage{hypcap}% it must be loaded after hyperref.
% Set up styles of URL: it should be placed after hyperref.
\urlstyle{same}

\addto\captionsitalian{\renewcommand{\contentsname}{ECONOMIA:}}

\usepackage{sphinxmessages}
\setcounter{tocdepth}{1}



\title{Appunti di Diritto ed Economia}
\date{05 mar 2026}
\release{3}
\author{Alessio Rubino}
\newcommand{\sphinxlogo}{\vbox{}}
\renewcommand{\releasename}{Release}
\makeindex
\begin{document}

\ifdefined\shorthandoff
  \ifnum\catcode`\=\string=\active\shorthandoff{=}\fi
  \ifnum\catcode`\"=\active\shorthandoff{"}\fi
\fi

\pagestyle{empty}
\sphinxmaketitle
\pagestyle{plain}
\sphinxtableofcontents
\pagestyle{normal}
\phantomsection\label{\detokenize{index::doc}}


\begin{figure}[htbp]
\centering

\noindent\sphinxincludegraphics[scale=0.6]{{diritto}.jpg}
\end{figure}

\sphinxstepscope


\chapter{1. La Moneta e il Baratto}
\label{\detokenize{1.La Moneta:la-moneta-e-il-baratto}}\label{\detokenize{1.La Moneta::doc}}
\begin{figure}[htbp]
\centering

\noindent\sphinxincludegraphics[scale=0.4]{{barattomoneta}.jpeg}
\end{figure}

\sphinxAtStartPar
Qualunque tipo di mercato c’è sempre un \sphinxstylestrong{venditore} e un \sphinxstylestrong{acquirente:}

\sphinxAtStartPar
il primo cerca di ottenere il massimo dal suo bene,
il secondo cerca di pagarlo il meno possibile.

\sphinxAtStartPar
Originariamente non esisteva il pagamento perché non c’era la moneta e quindi l’unica possibilità
di commercio è il baratto \sphinxstyleemphasis{(scambio di oggetti)}.

\sphinxAtStartPar
Il \sphinxstylestrong{baratto} ha però dei limiti perché può essere attuato con beni di identico valore e inoltre
perché alcuni beni erano difficilmente trasportabili.

\sphinxAtStartPar
Inizialmente è stato superato producendo beni che piacevano a tutti e che erano accettati nello
scambio di beni.

\sphinxAtStartPar
Successivamente questi beni sono stati sostituiti da pietre e metalli preziosi.

\sphinxAtStartPar
Ancora dopo i metalli sono state create delle monete, ad arrivare ai tempi moderni con l’uso (senza) metalli
preziosi e il valore è deciso dallo stato che lo emette.

\sphinxAtStartPar
Per questa ragione il cittadino deve accettare il pagamento dei beni con la moneta che ha emesso lo stato.

\sphinxAtStartPar
Ogni stato ha un suo livello economico che può essere più o meno elevato:
ci sono stati progrediti con un livello molto alto e stati, invece, chiamati in via di
sviluppo con un livello economico basso.


\section{1.1 Circolazione della Moneta}
\label{\detokenize{1.La Moneta:circolazione-della-moneta}}
\sphinxAtStartPar
Come tutti i beni presenti in uno Stato, anche la moneta ha un suo mercato dove ci sono soggetti che la acquistano
e la danno.

\sphinxAtStartPar
In generale se ci sono molte transazioni in un sistema, la moneta avrà una velocità di circolazione elevata.

\sphinxAtStartPar
In generale una velocità elevata è tipica dei sistemi espansivi in cui l’economia sta bene e vuole crescere,
invece, la politica restrittiva stagna nei mercati con problemi e con crescita ridotta.

\sphinxAtStartPar
La moneta in circolazione viene decisa dal Governo, mentre la sua circolazione dipende dagli intermediari
finanziari \sphinxstyleemphasis{(coloro che prendono la moneta e la utilizzano)}.

\sphinxAtStartPar
I principali intermediari sono:

\sphinxAtStartPar
\sphinxstylestrong{Le BANCHE}

\sphinxAtStartPar
\sphinxstylestrong{Le FINANZIARIE}

\sphinxAtStartPar
\sphinxstylestrong{I FONDI DI INVESTIMENTO}

\sphinxstepscope


\chapter{2. DIRITTI ECONOMICI}
\label{\detokenize{2.DIRITTI ECONOMICI-LAVORO:diritti-economici}}\label{\detokenize{2.DIRITTI ECONOMICI-LAVORO::doc}}

\section{2.1 \sphinxhyphen{} IL LAVORO}
\label{\detokenize{2.DIRITTI ECONOMICI-LAVORO:il-lavoro}}
\sphinxAtStartPar
PER POTER LAVORARE UN SOGGETTO DEVE ESSERE \sphinxstylestrong{ASSUNTO}.
\begin{description}
\sphinxlineitem{L’ASSUNZIONE AVVIENE ATTRAVERSO 2 CANALI:}\begin{enumerate}
\sphinxsetlistlabels{\arabic}{enumi}{enumii}{}{.}%
\item {} 
\sphinxAtStartPar
ISCRIZIONE NELLE \sphinxstylestrong{LISTE DI COLLOCAMENTO} nelle quali ci si può iscrivere dopo i 16 anni, con il consenso dei genitori, indicando i titoli posseduti e il lavoro che si vorrebbe svolgere.

\item {} 
\sphinxAtStartPar
\sphinxstylestrong{CHIAMATA DIRETTA}, quando il lavoratore prende accordi direttamente con il datore di lavoro. In questo caso è necessario il curriculum vitae in cui sono indicati tutti i titoli e le attività svolte.

\end{enumerate}

\end{description}

\sphinxAtStartPar
Se il lavoratore viene assunto viene fatto un \sphinxstylestrong{CONTRATTO DI LAVORO} in cui si deve riportare:
\begin{enumerate}
\sphinxsetlistlabels{\arabic}{enumi}{enumii}{}{.}%
\item {} 
\sphinxAtStartPar
IL NOME DEL LAVORATORE;

\item {} 
\sphinxAtStartPar
IL NOME DEL DATORE DI LAVORE;

\item {} 
\sphinxAtStartPar
LE MANSIONI CHE VENGONO ATTRIBUTE;

\item {} 
\sphinxAtStartPar
L’ORARIO DI LAVORO;

\item {} 
\sphinxAtStartPar
SE È UN CONTRATTO A TEMPO DETERMINATO (DATA INIZIO/FINE) O LAVORO INDETERMINATO;

\item {} 
\sphinxAtStartPar
DATA E FIRME.

\end{enumerate}


\section{2.2 \sphinxhyphen{} CCNL \sphinxhyphen{} Contratto Collettivo Nazionale di Lavoro}
\label{\detokenize{2.DIRITTI ECONOMICI-LAVORO:ccnl-contratto-collettivo-nazionale-di-lavoro}}
\sphinxAtStartPar
Tutti i DIRITTI e DOVERI dei lavoratori sono riportati in un contratto Generale che viene stipulato tra i
\sphinxstyleemphasis{Rappresentanti dei Lavori} e i \sphinxstyleemphasis{Rappresentanti dei Datori di Lavoro}.

\sphinxAtStartPar
I diritti sono INVIOLABILI e ogni accordo che NEGA questo è illegale.


\section{2.3 \sphinxhyphen{} I Sindacati}
\label{\detokenize{2.DIRITTI ECONOMICI-LAVORO:i-sindacati}}
\sphinxAtStartPar
A partire dagli anni 70” in Italia, nascono i sindacati che sono delle associazioni non riconosciute che hanno come scopo \sphinxstyleemphasis{rappresentare e tutetalare i lavoratori}.

\sphinxAtStartPar
Ogni lavoratore ha il DIRITTO di iscriversi ad un sindacato.

\sphinxAtStartPar
I principali sindacati in Italia sono 3 e sono chiaati \sphinxstyleemphasis{confederati} perchè hanno accordi tra di loro.
\begin{enumerate}
\sphinxsetlistlabels{\arabic}{enumi}{enumii}{}{.}%
\item {} 
\sphinxAtStartPar
CGIL:\sphinxhyphen{} di sinistra

\item {} 
\sphinxAtStartPar
CISL:\sphinxhyphen{} di centro

\item {} 
\sphinxAtStartPar
UIL :\sphinxhyphen{} di sinistra

\end{enumerate}

\sphinxAtStartPar
Il più importantebè la CGIL che ha più iscritti.
Negli anni sono nati altri sindacati definigti \sphinxstyleemphasis{autonomi} e negli ultimi anni anche i COBAS che sono sindacai di \sphinxstyleemphasis{BASE}.
Sono quelli più agguerriti perchè hanno pochi isciritti.


\section{2.4 \sphinxhyphen{} I DIRITTI DEI LAVORATORI}
\label{\detokenize{2.DIRITTI ECONOMICI-LAVORO:i-diritti-dei-lavoratori}}\begin{enumerate}
\sphinxsetlistlabels{\arabic}{enumi}{enumii}{}{.}%
\item {} 
\sphinxAtStartPar
DIRITTO di svolgere le «MANSIONI» scritte nel contratto;

\item {} 
\sphinxAtStartPar
DIRITTO di lavorare nell’ORARIO stabilito;

\item {} 
\sphinxAtStartPar
DIRITTO di alla RETRIBUZIONE;

\item {} 
\sphinxAtStartPar
DIRITTO al RIPOSO SETTIMANALE;

\item {} 
\sphinxAtStartPar
DIRITTO alle FERIE;

\item {} 
\sphinxAtStartPar
DIRITTO agli EVENTUALI STRAORDINARI;

\item {} 
\sphinxAtStartPar
DIRITTO alla SICUREZZA sul Lavoro;

\item {} 
\sphinxAtStartPar
DIRITTO di ASSENTARSI per «MALATTIA»;

\item {} 
\sphinxAtStartPar
DIRITTO delle DONNE ad ASSENTARSI quando nasce un figlio per un periodo PRIMA e DOPO il PARTO;

\item {} 
\sphinxAtStartPar
DIRITTO di entrambi i genitori di Assentarsi per assistere il FIGLIO APPENA NATO;

\item {} 
\sphinxAtStartPar
DIRITTO alla PENSIONE.

\end{enumerate}


\section{2.5 \sphinxhyphen{} DOVERI}
\label{\detokenize{2.DIRITTI ECONOMICI-LAVORO:doveri}}\begin{enumerate}
\sphinxsetlistlabels{\arabic}{enumi}{enumii}{}{.}%
\item {} 
\sphinxAtStartPar
DOVERE DI PRESTARE L’ATTIVITÀ in modo da rispettare le mansioni e l’orario stabilito

\item {} 
\sphinxAtStartPar
DOVERE DI RISPETTARE la sicurezza sul lavoro

\item {} 
\sphinxAtStartPar
DOVERE DI RETRIBUIRE IL DIPENDENTE

\item {} 
\sphinxAtStartPar
xxxxxxxxxx

\item {} 
\sphinxAtStartPar
POTERE DI DARE DISPOSIZIONI AI LAVORATORI

\item {} 
\sphinxAtStartPar
POTERE DI CONTROLLO PER VERIFICARE SE LE DISPOSIZIONI SONO STATE RISPETTATE

\item {} 
\sphinxAtStartPar
POTERE DISCIPLINARE, se il lavoratore NON rispetta le direttive, possono essere emesse delle sanzioni che variano a seconda del tipo di violazione fino al massimo che è il LICENZIAMENTO.

\end{enumerate}


\section{2.6 \sphinxhyphen{} CESSAZIONE DEL LAVORO}
\label{\detokenize{2.DIRITTI ECONOMICI-LAVORO:cessazione-del-lavoro}}
\sphinxAtStartPar
\sphinxstylestrong{LA CESSAZIONE DEL RAPPORTO DI LAVORO AVVIENE IN 5 CASI}:
\begin{enumerate}
\sphinxsetlistlabels{\arabic}{enumi}{enumii}{}{.}%
\item {} 
\sphinxAtStartPar
LICENZIAMENTO PER VARI MOTIVI

\item {} 
\sphinxAtStartPar
DIMISSIONI DEL LAVORATORE

\item {} 
\sphinxAtStartPar
MORTE DEL LAVORATORE

\item {} 
\sphinxAtStartPar
RAGGIUNGIMENTO DELL’ETÀ DELLA PENSIONE

\item {} 
\sphinxAtStartPar
SCADENZA DEL CONTRATTO A TEMPO DETERMINATO

\end{enumerate}

\sphinxstepscope


\chapter{3. DIRITTI DEGLI IMPRENDITORI}
\label{\detokenize{3.DIRITTI DEGLI IMPRENDITORI:diritti-degli-imprenditori}}\label{\detokenize{3.DIRITTI DEGLI IMPRENDITORI::doc}}
\sphinxAtStartPar
Nell’Art. 42 della Costituzione è prevista la tutela della \sphinxstylestrong{PROPRIETÀ PRIVATA} sia del singolo
che degli imprenditori.

\sphinxAtStartPar
Di conseguenza ognuno può fare quello che vuole con i propri beni senza che lo Stato intervenga.
C’è una sua protezione che lo Stato interviene sulla proprietà privata solo per scopo di
\sphinxstylestrong{PUBBLICA UTILITÀ}.

\sphinxAtStartPar
Così se c’è un interesse pubblico (costruire una strada) può togliere la proprietà al privato senza
che questi si possa opporre.

\sphinxAtStartPar
Al privato spetta un \sphinxstylestrong{INDENNIZZO} il cui importo è molto inferiore al valore reale del bene.

\sphinxAtStartPar
Questo procedimento viene attuato perché l’interesse collettivo prevale su quello privato.

\sphinxAtStartPar
Inoltre lo Stato riconosce all’imprenditore il \sphinxstylestrong{diritto di iniziare, cessare, qualunque tipo di attività}
senza interferenza dello Stato.

\sphinxAtStartPar
Solo in alcuni casi lo Stato vieta un’attività economica:
\begin{enumerate}
\sphinxsetlistlabels{\arabic}{enumi}{enumii}{}{.}%
\item {} 
\sphinxAtStartPar
VA CONTRO LA SALUTE

\item {} 
\sphinxAtStartPar
L’AMBIENTE

\item {} 
\sphinxAtStartPar
VA CONTRO LA DIGNITA” DELLA PERSONA

\item {} 
\sphinxAtStartPar
VA CONTRO LA LIBERTA” DELLA PERSONA

\end{enumerate}

\sphinxstepscope


\chapter{4. L’INFLAZIONE}
\label{\detokenize{4.L'INFLAZIONE:l-inflazione}}\label{\detokenize{4.L'INFLAZIONE::doc}}
\begin{figure}[htbp]
\centering

\noindent\sphinxincludegraphics[scale=1.0]{{inflazione}.jpeg}
\end{figure}

\sphinxAtStartPar
Ogni economia ha una sua inflazione, che è un meccanismo silenzioso e graduale:
una moneta perde di valore e diminuisce il suo potere d’acquisto.

\sphinxAtStartPar
Per questa ragione una moneta ha 2 tipi di valore:

\sphinxAtStartPar
\sphinxstylestrong{quello nominale} (che è quello espresso nella moneta)

\sphinxAtStartPar
e \sphinxstylestrong{quello reale} (il potere d’acquisto che corrisponde al valore reale).

\sphinxAtStartPar
Il potere d’acquisto che corrisponde al vaolre reale è strettamente legato a all’aumento dei prezzi.

\sphinxAtStartPar
Più aumentano i prezzi e meno vale la moneta.

\sphinxAtStartPar
L’inflazione può essere di più tipi :
\begin{enumerate}
\sphinxsetlistlabels{\arabic}{enumi}{enumii}{}{.}%
\item {} 
\sphinxAtStartPar
\sphinxstylestrong{STRISCAINATE} quando si parla di percentuali basse di inflazione;

\item {} 
\sphinxAtStartPar
\sphinxstylestrong{GALOPPANTE} quando di parla di decine di punti percentuali;

\item {} 
\sphinxAtStartPar
\sphinxstylestrong{IPER INFALZIONE} quando l’aumento è superiore al 50\% per un mese.

\end{enumerate}

\sphinxAtStartPar
L’inflazione può essere causata da fattori:

\sphinxAtStartPar
\sphinxstylestrong{Interni}: Aumento prezzi, salari, moneta in circolazione; sono le cause che possono aumentare
l’inflazione e lo Stato deve saper controllare.

\sphinxAtStartPar
\sphinxstylestrong{Esterni}: Cause che vengono da altri stati (es. guerre…); rispetto ai quali il singolo
stato non può fare assolutamente niente.


\section{4.1 \sphinxhyphen{} CAUSE INFLAZIONE}
\label{\detokenize{4.L'INFLAZIONE:cause-inflazione}}
\sphinxAtStartPar
In generale l’inflazione aumenta quando salgono i prezzi, e ciò avviene nei seguenti casi:

\sphinxAtStartPar
\sphinxstylestrong{Aumento dei salari}
\sphinxstylestrong{Aumento della moneta in circolazione}
\sphinxstylestrong{L’inflazione importata}, cioè da costi sper materie energetiche


\section{4.2 \sphinxhyphen{} MISURAZIONE INFLAZIONE}
\label{\detokenize{4.L'INFLAZIONE:misurazione-inflazione}}
\sphinxAtStartPar
Per stabilire il tasso di inflazione, c’è un organismo chiamato \sphinxstylestrong{ISTAT} in Italia che tiene sotto
controllo il mercato e le variazioni di prezzo.
In pratica ha preso come campione il prezzo di determinati beni che sono quelli che si consumano maggiormente
in Italia, chiamati \sphinxstylestrong{paniere economico}.

\begin{figure}[htbp]
\centering

\noindent\sphinxincludegraphics[scale=1.0]{{paniere}.jpeg}
\end{figure}

\sphinxAtStartPar
Periodicamente, ogni mese, l’ISTAT rileva la variazione di prezzo dei beni del paniere,
la variazione in Italia, fa la media e comunica la variazione \% dei prezzi.

\sphinxAtStartPar
In questo modo si ha la variazione ISTAT di ogni mese e allo quello annuale rispetto all’anno precedente.

\sphinxstepscope


\chapter{DOVERI COSTITUZIONALI}
\label{\detokenize{Doveri Costituzionali:doveri-costituzionali}}\label{\detokenize{Doveri Costituzionali::doc}}
\begin{figure}[htbp]
\centering

\noindent\sphinxincludegraphics[scale=0.6]{{dovericostituzionali}.jpeg}
\end{figure}

\sphinxAtStartPar
Ogni cittadino italiano ha dei doveri.

\sphinxAtStartPar
Il \sphinxstylestrong{PRIMO} di questi è quello di \sphinxstylestrong{contribuire alle spese dello Stato}
attraverso i tributi (imposte, tasse e contributi).

\sphinxAtStartPar
Quest’obbligo è previsto \sphinxstylestrong{nell’articolo 53 della Costituzione} e deve rispettare 2 principi:
\begin{enumerate}
\sphinxsetlistlabels{\arabic}{enumi}{enumii}{}{.}%
\item {} 
\sphinxAtStartPar
\sphinxstylestrong{CAPACITÀ CONTRIBUTIVA}: la contribuzione deve avvenire in base al reddito e quindi i più ricchi contribuiscono di più;

\item {} 
\sphinxAtStartPar
\sphinxstylestrong{PROGRESSIVITÀ} nel senso che le aliquote, la \% che si deve versare allo Stato, cresce con l’aumentare del reddito.

\end{enumerate}

\sphinxAtStartPar
Il \sphinxstylestrong{SECONDO} DOVERE: \sphinxstylestrong{difesa della Patria}. In base al quale ogni cittadino deve difendere
la Patria in caso d’attacco da un altro Stato.
Fino agli anni 2000 tutti i cittadini maschi dovevano prestare servizio militare per 1 anno dal compimento dei 18 anni. Poi quest’obbligo viene cancellato.

\sphinxAtStartPar
Attualmente l’esercito italiano è composto soltanto da volontari che vengono retribuiti dallo Stato. Però in caso di guerra ogni cittadino maggiorenne può essere chiamato a combattere. In caso di rifiuto c’è il \sphinxstylestrong{reato di diserzione}, cioè l’ergastolo. Se si verifica durante il combattimento c’è la pena di morte.

\sphinxAtStartPar
Il \sphinxstylestrong{TERZO} DOVERE: \sphinxstylestrong{fedeltà alla Repubblica} consiste nell’obbligo di ogni cittadino di non compiere attività che possa danneggiare il proprio Stato. Rientra in quest’obbligo quello di rispettare le leggi perché in questo modo si garantisce il buon funzionamento dello Stato.
Fammi sapere se hai bisogno di altro!

\sphinxstepscope


\chapter{Il Sistema ELETTORALE}
\label{\detokenize{Sistema Elettorale:il-sistema-elettorale}}\label{\detokenize{Sistema Elettorale::doc}}
\sphinxAtStartPar
Il sistema elettorale è il modo in cui si vota per eleggere i parlamentari ed è diviso nelle seguenti forme:

\sphinxAtStartPar
\sphinxstylestrong{MAGGIORITARIO}

\sphinxAtStartPar
È manipolato e in Italia viene diviso in:

\sphinxAtStartPar
\sphinxstylestrong{Uninominale}: quando in un collegio vince solo quello che ha ricevuto la maggioranza dei voti.

\sphinxAtStartPar
\sphinxstylestrong{Plurinominale}: quando nel collegio non vengono eletti solo i rappresentanti della maggioranza,
ma è democratico perché anche l’opposizione ne ha diritto.
Questo sistema favorisce la stabilità del Governo, ma ha il difetto di non riflettere la volontà dei partiti.

\sphinxAtStartPar
\sphinxstylestrong{PROPORZIONALE}

\sphinxAtStartPar
In base al quale i parlamentari vengono ripartiti in proporzione ai voti ottenuti dai partiti.
Questo sistema è più democratico in quanto rispetta i voti, ma genera instabilità del Governo
perché perdono tutti i partiti e lo Stato perde credibilità perché tutti parlano male.

\sphinxAtStartPar
Permette la rappresentanza anche ai partiti più piccoli.

\sphinxAtStartPar
\sphinxstylestrong{MISTO/MATTARELLA}

\sphinxAtStartPar
In cui l’elezione del 75\% dei parlamentari avviene con l’uninominale (cioè il maggioritario) e il 25\% con il
proporzionale.

\sphinxstepscope


\chapter{DIRITTO DI VOTO ED ELEZIONI}
\label{\detokenize{Diritto di voto ed elezioni:diritto-di-voto-ed-elezioni}}\label{\detokenize{Diritto di voto ed elezioni::doc}}
\begin{figure}[htbp]
\centering

\noindent\sphinxincludegraphics[scale=0.6]{{dirittodivoto}.jpeg}
\end{figure}

\sphinxAtStartPar
L’Italia è una Repubblica democratica nella quale si stabilisce il potere e dalla Costituzione
che la sovranità è esercitata attraverso il voto, ovvero la democrazia.

\sphinxAtStartPar
Il voto ha 4 caratteristiche:

\sphinxAtStartPar
\sphinxstylestrong{UNIVERSALE}, tutti gli uomini maggiorenni e capaci sono chiamati a votare;

\sphinxAtStartPar
\sphinxstylestrong{UGUALE}, nel senso che il voto del ricco è uguale a quello del povero;

\sphinxAtStartPar
\sphinxstylestrong{PERSONALE}, cioè ciascuno deve votare personalmente e nessuno deve delegare;

\sphinxAtStartPar
\sphinxstylestrong{LIBERO e SEGRETO}, nel senso che ognuno vota chi vuole e non è tenuto a dirlo;

\sphinxAtStartPar
Come cittadino si ha il diritto al voto e anche un \sphinxstylestrong{dovere civico} di votare.

\sphinxstepscope


\chapter{IL PARLAMENTO}
\label{\detokenize{IL PARLAMENTO:il-parlamento}}\label{\detokenize{IL PARLAMENTO::doc}}
\begin{figure}[htbp]
\centering

\noindent\sphinxincludegraphics[scale=0.6]{{parlamento}.jpeg}
\end{figure}

\sphinxAtStartPar
Il \sphinxstylestrong{Parlamento} è il titolare del potere legislativo.
In Italia è composto da due camere che hanno lo stesso potere, cioè bicameralismo perfetto, e sono la:
\begin{itemize}
\item {} 
\sphinxAtStartPar
\sphinxstylestrong{CAMERA DEI DEPUTATI}, sede Montecitorio

\item {} 
\sphinxAtStartPar
\sphinxstylestrong{SENATO DELLA REPUBBLICA}, sede Palazzo Madama

\end{itemize}

\sphinxAtStartPar
Normalmente le 2 camere si riuniscono e lavorano in modo diviso.
Solo in alcuni casi si riuniscono in seduta comune, ad esempio quando:
\begin{enumerate}
\sphinxsetlistlabels{\arabic}{enumi}{enumii}{}{.}%
\item {} 
\sphinxAtStartPar
NOMINARE PDR;

\item {} 
\sphinxAtStartPar
ELEGGERE un 1/3 ELEMENTO DELLA CORTE COSTITUZIONALE;

\item {} 
\sphinxAtStartPar
QUANDO deve eleggere i membri del CSM (Consiglio Superiore della Magistratura);

\item {} 
\sphinxAtStartPar
QUANDO deve mettere in ACCUSA IL PRESIDENTE DELLA REPUBBLICA.

\end{enumerate}

\sphinxAtStartPar
Il Parlamento ha una carica di 5 anni (chiamata legislatura) e si possono rieleggere tutti i giorni. Nella Camera che è composta da 400 deputati, sono tutti eletti dal popolo, mentre il Senato, composto da 200 membri, i 5 Senatori a vita e gli ex Presidenti della Repubblica non sono eletti dal popolo e entrano in carica per tutta la vita.


\section{FUNZIONAMENTO DELLE CAMERE}
\label{\detokenize{IL PARLAMENTO:funzionamento-delle-camere}}
\sphinxAtStartPar
Ogni camera ha un \sphinxstylestrong{Presidente} che ha il compito di dirigere i lavori, dare la parola e togliere la parola
durante le discussioni, è eletto a maggioranza dalla camera di appartenenza.

\begin{figure}[htbp]
\centering

\noindent\sphinxincludegraphics[scale=1.0]{{presidenti}.jpeg}
\end{figure}

\sphinxAtStartPar
Successivamente ci sono i \sphinxstylestrong{gruppi parlamentari} e c’è l’unione di più parlamentari in un tema politico.
Ogni gruppo decide ogni rappresentante o li rappresenta. I gruppi deliberano in piena libertà politica durante
le votazioni ma negli orari parlamentari non sono vincolati al rispetto delle indicazioni del gruppo.

\begin{figure}[htbp]
\centering

\noindent\sphinxincludegraphics[scale=0.3]{{gruppiparlmentari}.jpeg}
\end{figure}

\sphinxAtStartPar
Tra le \sphinxstylestrong{commissioni parlamentari}, gruppi di deputati e senatori esaminano una materia specifica che servono
da filtro tra le proposte di legge e l’esame sulle stesse.
Le commissioni convocate per materia ritengono che la proposta sia scadibile e viene poi la prova del che
l’ordinamento diviene stretto.

\sphinxAtStartPar
Esistono 2 tipi di commissioni:

\sphinxAtStartPar
\sphinxstylestrong{PERMANENTI}, che fanno il filtro;

\sphinxAtStartPar
\sphinxstylestrong{INCHIESTA}, quando ci sono problemi a livello nazionale in cui bisogna indagare.

\begin{figure}[htbp]
\centering

\noindent\sphinxincludegraphics[scale=1.0]{{commissioni}.jpeg}
\end{figure}

\sphinxAtStartPar
\sphinxstylestrong{IMMUNITA” Parlamentare}

\sphinxAtStartPar
Il parlamentare quando esprime il suo pensiero è insindacabile, quindi nessuno lo può processare per quello che dice.
Inoltre se commette dei reati e deve essere processato, occorre ottenere l’autorizzazione della camera di
appartenenza se non viene colto in flagranza o in atto.

\begin{figure}[htbp]
\centering

\noindent\sphinxincludegraphics[scale=1.0]{{immunita}.jpg}
\end{figure}

\sphinxAtStartPar
Infine ogni parlamnetare prendo un indennità (Stipendio) il cui importo è deciso dalla camera di appartenenza.

\sphinxstepscope


\chapter{PROCEDIMENTO LEGISLATIVO (LEGGE)}
\label{\detokenize{PROCEDIMENTO LEGISLATIVO:procedimento-legislativo-legge}}\label{\detokenize{PROCEDIMENTO LEGISLATIVO::doc}}
\sphinxAtStartPar
Per approvare una legge il parlamento segue una procedura chiamata \sphinxstylestrong{«iter legis»} che ha diverse fasi:

\sphinxAtStartPar
\sphinxstylestrong{1ª FASE = INIZIATIVA}, ci sono 5 soggetti che possono fare una proposta di legge e sono:
\begin{itemize}
\item {} 
\sphinxAtStartPar
\sphinxstylestrong{GOVERNO}, con i disegni di legge;

\item {} 
\sphinxAtStartPar
\sphinxstylestrong{ogni parlamentare};

\item {} 
\sphinxAtStartPar
\sphinxstylestrong{le regioni}, per le materie sui cui hanno competenza;

\item {} 
\sphinxAtStartPar
\sphinxstylestrong{cittadini elettori}, 50.000;

\item {} 
\sphinxAtStartPar
\sphinxstylestrong{CNEL}.

\end{itemize}

\sphinxAtStartPar
\sphinxstylestrong{2ª FASE}, le proposte passano ad una delle due camere.

\sphinxAtStartPar
\sphinxstylestrong{3ª FASE}, la camera che riceve la proposta decide di esaminarla all’interno della commissione parlamentare della materia che la discute e alla fine decide se farla andare avanti o bloccarla;

\sphinxAtStartPar
\sphinxstylestrong{4ª FASE}, se la commissione decide di esaminarla, viene convocato un avviso, cioè stabilisce il giorno del suo esame;

\sphinxAtStartPar
\sphinxstylestrong{5ª FASE}, la discussione è proclamata spesso alla conferenza dei capi gruppo;

\sphinxAtStartPar
\sphinxstylestrong{6ª FASE}, nel giorno stabilito inizia la discussione. Organi parlamentari possono chiedergli di intervenire oppure può presentare delle richieste di modifica chiamate «emendamenti»;

\sphinxAtStartPar
\sphinxstylestrong{7ª FASE}, la proposta viene discusso e votata.

\sphinxAtStartPar
Certo, ecco la trascrizione degli appunti che mi hai fornito nell’immagine:

\sphinxAtStartPar
\sphinxstylestrong{8° FASE:}
* discussione gli eventuali emendamenti;
* si passa alla votazione che si effettua in questo modo:
* si votano gli emendamenti;
* si votano i singoli articoli;
* si vota tutta la legge nel suo complesso.

\sphinxAtStartPar
\sphinxstylestrong{9° FASE:}
* il parlamento approva la legge vota in altro ramo dove viene ripetuto lo stesso procedimento;

\sphinxAtStartPar
\sphinxstylestrong{10° FASE:}
* se l’altro ramo approva degli emendamenti la legge deve ritornare all’altro ramo fino a quando entrambi non approvano lo stesso testo. Il tutto deve avvenire in un tempo massimo che viene stabilito in ogni caso, passato il quale il provvedimento decade.

\sphinxAtStartPar
\sphinxstylestrong{11° FASE:}
* se la legge è approvata su tutti e due rami passa dal PdR che esamina per accettare se ci sono contrasti con la Costituzione e ci possono essere 2 motivi:
* non ci sono contrasti e promulga la legge;
* ci sono contrasti e in questo caso viene mandato al parlamento con un messaggio dove è possibile il contrasto.

\sphinxAtStartPar
\sphinxstylestrong{Il Parlamento ha 3 possibilità:}
* Approva la legge con le modifiche indicate dal PdR che quindi la promulga.
* Il Parlamento riapprova la legge con lo stesso testo e in questo caso il PdR può fare 2 cose:
* Promulga la legge;
* SI RIFIUTA E SI DEVE DIMETTERE.

\sphinxAtStartPar
\sphinxstylestrong{12° FASE} SE LA LEGGE È APPROVATA DA TUTTI E 2 (i rami del Parlamento), VIENE PASSATO
(al Presidente della Repubblica) CHE LA ESAMINA PER ACCETTARE CHE NON CI SIANO CONTRASTI CON LA COSTITUZIONE.
Ci sono 2 ipotesi:
\sphinxhyphen{} NON CI SONO CONTRASTI E {[}il PDR la promulga{]}
\sphinxhyphen{} CI SONO CONTRASTI e in questo caso rinvia al Parlamento con la motivazione del perché {[}ci sono{]} contrasti.
Il Parlamento ha 3 possibilità:
\sphinxhyphen{} RIAPPROVA la legge con le modifiche e la promulga
\sphinxhyphen{} IL PARLAMENTO RIAPPROVA la legge con lo stesso testo. In questo caso il PDR può fare 2 cose:
1. PROMULGA la legge
2. SI RIFIUTA E SI DEVE DIMETTERE

\sphinxAtStartPar
A questo punto, se la legge è accettata, viene inviata ad una \sphinxstylestrong{GAZZETTA UFFICIALE} che la pubblica e dopo \sphinxstylestrong{45 GIORNI} viene pubblicata. Questo periodo si chiama {[}Vacatio Legis{]} e serve ai cittadini per \sphinxstylestrong{DA CONOSCERE LA LEGGE}. In caso di urgenza, la legge può entrare in vigore anche lo stesso giorno della pubblicazione.

\sphinxstepscope


\chapter{ITER LEGISLATIVO ABBREVIATO E AGGRAVATO}
\label{\detokenize{ITER LEGISLATIVO ABBREVIATO E AGGRAVATO:iter-legislativo-abbreviato-e-aggravato}}\label{\detokenize{ITER LEGISLATIVO ABBREVIATO E AGGRAVATO::doc}}
\sphinxAtStartPar
OLTRE ALL’ITER LEGIS ORDINARIO, la nostra Costituzione ne prevede altri 2 che sono quello \sphinxstylestrong{ABBREVIATO}
e quello \sphinxstylestrong{AGGRAVATO}.

\sphinxAtStartPar
Il primo (Abbreviato) si utilizza in situazioni di urgenza, quando non è possibile adottare un provvedimento con
iter ordinario perché è troppo {[}lento{]}.

\sphinxAtStartPar
In questa ipotesi il Governo può adottare un \sphinxstylestrong{DECRETO LEGGE} che ha gli effetti di una legge ma non
ha la forza della legge perché non è adottato dal Parlamento.

\sphinxAtStartPar
Questo decreto ha un’efficacia limitata nel tempo di \sphinxstylestrong{60 GIORNI}.
Entro questo periodo il Parlamento lo deve convertire in legge e se non lo fa perde gli effetti.

\sphinxAtStartPar
Il procedimento \sphinxstylestrong{AGGRAVATO} è quello più complesso e serve per approvare le \sphinxstylestrong{LEGGI DI REVISIONE COSTITUZIONALE}.

\sphinxAtStartPar
Il procedimento è il seguente:
\begin{enumerate}
\sphinxsetlistlabels{\arabic}{enumi}{enumii}{}{.}%
\item {} 
\sphinxAtStartPar
Si segue l’iter ordinario e quindi la legge deve essere approvata da entrambi i rami ed è sufficiente la maggioranza semplice, cioè la metà più uno.

\item {} 
\sphinxAtStartPar
Una volta approvata la legge non entra in vigore perché non {[}è definitiva{]}.

\item {} 
\sphinxAtStartPar
Bisogna attendere almeno \sphinxstylestrong{3 MESI}, dopodiché si deve rifare lo stesso procedimento e quindi l’approvazione di entrambi i rami, in pratica occorre una \sphinxstylestrong{DOPPIA APPROVAZIONE} da entrambi.

\item {} 
\sphinxAtStartPar
SE ANCHE NELLA SECONDA APPROVAZIONE la legge viene nuovamente approvata ci sono 2 possibilità:
\begin{quote}
\begin{itemize}
\item {} 
\sphinxAtStartPar
SE VIENE APPROVATA CON LA \sphinxstylestrong{MAGGIORANZA SEMPLICE} NON ENTRA IN VIGORE perché entro 3 mesi può essere richiesto un \sphinxstylestrong{REFERENDUM CONSERVATIVO}.

\end{itemize}

\sphinxAtStartPar
Possono richiedere il referendum i seguenti:
\begin{enumerate}
\sphinxsetlistlabels{\arabic}{enumii}{enumiii}{}{.}%
\item {} 
\sphinxAtStartPar
1/5 DEI {[}MEMBRI DI UNA CAMERA{]}

\item {} 
\sphinxAtStartPar
5 CONSIGLI REGIONALI

\item {} 
\sphinxAtStartPar
500.000 CITTADINI ELETTORI

\end{enumerate}
\end{quote}

\end{enumerate}
\begin{itemize}
\item {} 
\sphinxAtStartPar
SE VIENE CHIESTO IL REFERENDUM OCCORRE CHE ALMENO IL 50\% PIÙ 1 DEGLI {[}AVENTI DIRITTO{]} VANO A VOTARE E LA MAGGIORANZA VOTA LA LEGGE.

\item {} 
\sphinxAtStartPar
SE NON È RAGGIUNTA questa maggioranza la legge \sphinxstylestrong{NON ENTRA} {[}in vigore{]}.

\end{itemize}
\begin{enumerate}
\sphinxsetlistlabels{\arabic}{enumi}{enumii}{}{.}%
\setcounter{enumi}{4}
\item {} 
\sphinxAtStartPar
SE ENTRO {[}TRE MESI{]} IL TEMPO NESSUNO CHIEDE IL REFERENDUM LA LEGGE ENTRA {[}in vigore{]}.

\item {} 
\sphinxAtStartPar
SE NELLA 2° VOTAZIONE SI RAGGIUNGE LA \sphinxstylestrong{MAGGIORANZA QUALIFICATA} {[}2/3{]} ENTRA IN VIGORE {[}subito{]}.

\end{enumerate}

\sphinxstepscope


\chapter{IL REFERENDUM}
\label{\detokenize{IL REFERENDUM:il-referendum}}\label{\detokenize{IL REFERENDUM::doc}}
\sphinxAtStartPar
In Italia il popolo solo in pochissimi casi interviene direttamente sull’attività dello Stato. Il referendum è uno di questi casi. È uno strumento con il quale è possibile annullare una legge approvata dal Parlamento. Ci sono leggi che non possono essere annullate per referendum, e tra queste:
\begin{itemize}
\item {} 
\sphinxAtStartPar
\sphinxstylestrong{MATERIA TRIBUTARIA};

\item {} 
\sphinxAtStartPar
\sphinxstylestrong{MATERIA PENALE};

\item {} 
\sphinxAtStartPar
\sphinxstylestrong{DISCIPLINANO i RAPPORTI INTERNAZIONALI};

\item {} 
\sphinxAtStartPar
\sphinxstylestrong{LEGGI PER IL FUNZIONAMENTO DELLO STATO}.

\end{itemize}

\sphinxAtStartPar
Per chiedere un referendum è necessaria la richiesta firmata da \sphinxstylestrong{500.000 cittadini elettori}.
La richiesta con le firme viene trasmessa alla Cassazione che controlla la regolarità.

\sphinxAtStartPar
Se la Cassazione dice che è tutto apposto la richiesta passa alla Corte Costituzionale la quale accetta se è
ammissibile e quindi che non si tratti uno dei divieti, e se ritiene che è tutto apposto.

\sphinxAtStartPar
Il PDR fissa la data della votazione.

\sphinxAtStartPar
In quel giorno tutti i cittadini elettori possono recarsi a votare.
Perché un referendum sia valido occorre che vada a votare il 50\% + 1 dei cittadini elettori.

\sphinxAtStartPar
Se non è raggiunto il \sphinxstylestrong{QUORUM} il referendum è nullo, se invece è raggiunto bisogna vedere la maggioranza come abbia votato:
se la maggioranza dice “SÌ”, la legge è {[}approvata e cessa di esistere{]}.

\sphinxstepscope


\chapter{Il Governo: Procedura di Formazione}
\label{\detokenize{Il GOVERNO - Formazione:il-governo-procedura-di-formazione}}\label{\detokenize{Il GOVERNO - Formazione::doc}}
\begin{figure}[htbp]
\centering

\noindent\sphinxincludegraphics[scale=1.0]{{governo}.jpg}
\end{figure}

\sphinxAtStartPar
Per formare il \sphinxstylestrong{Governo} occorre nominare il \sphinxstylestrong{Presidente del Consiglio dei Ministri}, il quale sceglie i Ministri.

\sphinxAtStartPar
Per scegliere il P.d.C. (Presidente del Consiglio) c’è una procedura chiamata \sphinxstylestrong{Consultazione}.

\sphinxAtStartPar
Il P.d.R. (Presidente della Repubblica) convoca i rappresentanti dei partiti e gli chiede quale soggetto
possa essere nominato P.d.C., così da ogni partito indica le sue preferenze.

\sphinxAtStartPar
Il P.d.R. nomina quello che è appoggiato dalla maggioranza dei parlamentari.

\sphinxAtStartPar
Il P.d.R. chiama il soggetto e chiede se accetta l’incarico.

\sphinxAtStartPar
Se lo accetta, deve scegliere i Ministri.

\sphinxAtStartPar
Quando i partiti si sono accordati, il P.d.C. fa un elenco e lo porta al P.d.R., il quale può fare delle
osservazioni su alcuni nomi se li ritiene non adatti.

\sphinxAtStartPar
Se non fa obiezioni, il P.d.C. e i Ministri devono giurare fedeltà alla Costituzione davanti al P.d.R.

\sphinxAtStartPar
Successivamente, P.d.C. e Ministri devono comparire davanti alle Camere per ottenere la fiducia.

\sphinxAtStartPar
Ogni Camera vota e, se entrambe danno la fiducia, il Governo entra in funzione.

\sphinxstepscope


\chapter{DIRITTI COSTITUZIONALI}
\label{\detokenize{DIRITTI COSTITUZIONALI:diritti-costituzionali}}\label{\detokenize{DIRITTI COSTITUZIONALI::doc}}
\sphinxAtStartPar
Dopo la II guerra mondiale l’Italia ha scelto di essere una Repubblica e ha approvato una nuova COSTITUZIONE.

\sphinxAtStartPar
In questa Costituzione si è dovuto scegliere il modo di governare l’Italia.

\sphinxAtStartPar
SI poteva scegliere un sistema presidenziale in cui tutti i poteri sono del Presidente, oppure vice presidenziale come in Francia ma entrambi questi modelli sono stati scartati per paura di tornare a una dittatura Fascista in cui un solo soggetto aveva portato l’Italia alla rovina punto per tale ragione si è scelto un sistema parlamentare che è una forma di governo più Lento ma più sicuro che si fonda sui presupposti :

\sphinxAtStartPar
Separazione dei poteri: che compongono la sovranità sono stati affidati a tre organi diversi Parlamento governo e magistratura con l’aggiunta di due altri organi di controllo e sono il presidente della Repubblica e la Corte Costituzionale.

\sphinxAtStartPar
In Italia il sistema di governo è fondato su 5 organi fra loro interconnessi che con che controllano circolarmente con questo sistema l’organo centrale e il Parlamento.

\begin{figure}[htbp]
\centering

\noindent\sphinxincludegraphics[scale=0.5]{{PoteriCostituzionali}.jpg}
\end{figure}
\begin{enumerate}
\sphinxsetlistlabels{\arabic}{enumi}{enumii}{}{.}%
\item {} 
\sphinxAtStartPar
Sistema parlamentare in base al quale i poteri principali cioè li ha il Parlamento

\item {} 
\sphinxAtStartPar
rappresentanza politica in base al quale il Parlamento rappresenta i cittadini politicamente

\item {} 
\sphinxAtStartPar
partiti politici sono associazioni non riconosciute ciascuno dei quali rappresenta un tipo di ideologia

\end{enumerate}

\sphinxstepscope


\chapter{Il Governo}
\label{\detokenize{IL GOVERNO:il-governo}}\label{\detokenize{IL GOVERNO::doc}}
\sphinxAtStartPar
Per formarlo occorre nominare il Presidente del Consiglio dei Ministri, il quale a sua volta sceglie i ministri.

\sphinxAtStartPar
Per scegliere il Presidente del Consiglio bisogna seguire le consultazioni:

\sphinxAtStartPar
Il PdR (Presidente della Repubblica) convoca i rappresentanti dei partiti e gli chiede quale soggetto vorrebbero che
fosse nominato PdC.

\sphinxAtStartPar
Dopo che i partiti indicano le preferenze, il PdR chiama il soggetto che vuole la maggioranza dei parlamentari e
gli dà l’incarico.

\sphinxAtStartPar
Se lo accetta, deve decidere chi sono i ministri; quando ha l’elenco lo porta al PdR, il quale può fare
delle osservazioni se non li ritiene adatti.

\sphinxAtStartPar
Se il PdR non fa rilievi, il PdC e i ministri devono giurare davanti a lui.

\sphinxAtStartPar
Successivamente, PdC e ministri devono comparire davanti alle Camere per ottenere la fiducia.

\sphinxAtStartPar
Ogni Camera vota: se entrambe danno la fiducia, il governo entra in funzione, altrimenti no.

\sphinxAtStartPar
Se, quando il governo è in funzione, il Parlamento gli revoca la fiducia, il governo cade e cessa la
legislatura nel momento in cui il Parlamento ha tolto la fiducia che gli aveva dato.

\sphinxstepscope


\chapter{Il potere normativo del Governo}
\label{\detokenize{IL POTERE NORMATIVO DEL GOVERNO:il-potere-normativo-del-governo}}\label{\detokenize{IL POTERE NORMATIVO DEL GOVERNO::doc}}
\sphinxAtStartPar
Normalmente il Governo non ha potere legislativo perché appartiene al Parlamento.

\sphinxAtStartPar
Tuttavia, in casi specifici, il Governo può adottare due provvedimenti che hanno la forza (ma non la sostanza) della legge:
\begin{itemize}
\item {} 
\sphinxAtStartPar
\sphinxstylestrong{Decreto Legge}: Il Parlamento ha tempi molto lunghi per approvare le leggi.

\end{itemize}

\sphinxAtStartPar
In casi d’urgenza, il Governo può adottare un provvedimento che ha una durata di 60 giorni, entro la quale deve essere approvato (convertito in legge) dal Parlamento, altrimenti decade.
\begin{itemize}
\item {} 
\sphinxAtStartPar
\sphinxstylestrong{Decreto Legislativo}: Ci sono materie complesse che richiedono tempi molto lunghi per l’esame e l’approvazione.

\end{itemize}

\sphinxAtStartPar
Per non bloccare l’attività del Parlamento, esso può delegare il Governo per approvare la legge tramite una legge delega. La delega deve però contenere 3 condizioni:
\begin{enumerate}
\sphinxsetlistlabels{\arabic}{enumi}{enumii}{}{.}%
\item {} 
\sphinxAtStartPar
L’oggetto della legge (la materia specifica);

\item {} 
\sphinxAtStartPar
I criteri e i principi direttivi della legge;

\item {} 
\sphinxAtStartPar
Il termine massimo (scadenza della delega), entro il quale il Governo deve approvare il decreto.

\end{enumerate}

\sphinxstepscope


\chapter{POTERE REGOLAMENTARE DEL GOVERNO}
\label{\detokenize{POTERE REGOLAMENTARE DEL GOVERNO:potere-regolamentare-del-governo}}\label{\detokenize{POTERE REGOLAMENTARE DEL GOVERNO::doc}}
\sphinxAtStartPar
Può accadere che una legge sia complessa e quindi richieda per la sua concreta attuazione, dei provvedimenti a
ggiuntivi chiamati regolamenti.

\sphinxAtStartPar
Questi provvedimenti sono chiamati fonti secondarie del diritto perché sono sottoposte alla legge e quindi non possono disporre nulla in contrasto con essa.

\sphinxAtStartPar
I regolamenti sono adottati dal governo in 2 modi:

\sphinxAtStartPar
Da ogni singolo ministro, a seconda della materia e si chiamano \sphinxstylestrong{DM (decreto ministeriale)};

\sphinxAtStartPar
Da tutto il governo, si chiamano \sphinxstylestrong{DPCM (decreto presidente Consiglio dei ministri).}

\sphinxAtStartPar
Nei rapporti interni il DPCM prevale sul DM.

\sphinxAtStartPar
I regolamenti hanno delle caratteristiche e dei limiti:
\begin{enumerate}
\sphinxsetlistlabels{\arabic}{enumi}{enumii}{}{.}%
\item {} 
\sphinxAtStartPar
Non possono andare contro la legge;

\item {} 
\sphinxAtStartPar
Sono attuativi di una legge;

\item {} 
\sphinxAtStartPar
In alcuni casi possono integrare la legge senza modificarla.

\end{enumerate}

\sphinxAtStartPar
Per non bloccare l’attività del Parlamento, esso può delegare il Governo ad approvare la legge
mettendo però 3 condizioni:

\sphinxAtStartPar
\sphinxstylestrong{L’oggetto della legge;}

\sphinxAtStartPar
\sphinxstylestrong{I principi e i criteri della legge;}

\sphinxAtStartPar
\sphinxstylestrong{La scadenza della delega} e cioè il termine massimo oltre il quale il Governo non può approvare
la scadenza della legge.

\sphinxstepscope


\chapter{La crisi di governo}
\label{\detokenize{LA CRISI DI GOVERNO:la-crisi-di-governo}}\label{\detokenize{LA CRISI DI GOVERNO::doc}}
\sphinxAtStartPar
La crisi di governo si verifica in una situazione di contrasto tra il Parlamento e il Governo, oppure all’interno
della maggioranza stessa.

\sphinxAtStartPar
In questa ipotesi si verifica la rottura del rapporto di fiducia: non c’è più una maggioranza che lo sostiene.

\sphinxAtStartPar
Le crisi possono essere:
\begin{itemize}
\item {} 
\sphinxAtStartPar
\sphinxstylestrong{Parlamentare}: quando si verifica all’interno del Parlamento (es. voto di sfiducia).

\item {} 
\sphinxAtStartPar
\sphinxstylestrong{Extraparlamentare}: quando la crisi nasce al di fuori del Parlamento (es. dimissioni spontanee del Presidente del Consiglio per contrasti politici tra i partiti della coalizione).

\end{itemize}

\sphinxstepscope


\chapter{La Fiducia e gli strumenti parlamentari}
\label{\detokenize{La Fiducia e gli strumenti parlamentari:la-fiducia-e-gli-strumenti-parlamentari}}\label{\detokenize{La Fiducia e gli strumenti parlamentari::doc}}
\sphinxAtStartPar
Quando il governo si presenta alle Camere, espone il proprio programma.

\sphinxAtStartPar
Se il Parlamento ritiene che tale programma non venga rispettato, può utilizzare uno strumento chiamato mozione di
sfiducia.

\sphinxAtStartPar
\sphinxstylestrong{Mozione di sfiducia}: È un atto attraverso il quale un decimo dei parlamentari può chiedere che si discuta
della revoca della fiducia al Governo.
Si può presentare in qualunque momento e in qualunque Camera.
\sphinxstyleemphasis{Dopo la discussione si vota:}

\sphinxAtStartPar
Se la maggioranza \sphinxstylestrong{approva} la mozione, il governo \sphinxstylestrong{decade};
Altrimenti, il governo rimane in carica.

\sphinxAtStartPar
\sphinxstylestrong{La Questione di Fiducia}

\sphinxAtStartPar
Esiste un altro strumento attraverso il quale il Governo può \sphinxstyleemphasis{«sfidare»} il Parlamento, spesso usato per sbloccare
l’iter di una legge: la questione di fiducia.

\sphinxAtStartPar
\sphinxstylestrong{Ostruzionismo}: Durante il procedimento legislativo, può accadere che l’opposizione presenti \sphinxstylestrong{migliaia} di emendamenti (proposte di modifica) per bloccare l’approvazione di una legge.

\sphinxAtStartPar
\sphinxstylestrong{Applicazione della fiducia}: Per far decadere tutti gli emendamenti in un colpo solo, il Governo \sphinxstyleemphasis{«pone la questione di fiducia»}
sulla legge.

\sphinxAtStartPar
In pratica, il governo dichiara: \sphinxstyleemphasis{«O approvate la legge così com’è, oppure mi dimetto».}

\sphinxAtStartPar
Se il Parlamento vota a favore, la legge è approvata e gli emendamenti decadono; se vota contro, il Governo è obbligato a dimettersi.

\sphinxAtStartPar
Nota sintetica a margine (presente nel testo):
Il rapporto di fiducia è fondamentale:
\begin{itemize}
\item {} 
\sphinxAtStartPar
\sphinxstylestrong{Mozione di sfiducia}: Parte dal Parlamento contro il Governo.

\item {} 
\sphinxAtStartPar
\sphinxstylestrong{Questione di fiducia}: Parte dal Governo verso il Parlamento (spesso per superare l’ostruzionismo).

\end{itemize}

\sphinxstepscope


\chapter{Corte Costituzionale}
\label{\detokenize{Corte Costituzionale:corte-costituzionale}}\label{\detokenize{Corte Costituzionale::doc}}
\sphinxAtStartPar
Insieme al PdR, la Corte è un organo di controllo.
È composta da 15 giudici, nominati per:
\begin{itemize}
\item {} 
\sphinxAtStartPar
\sphinxstylestrong{5} dal PdR (Presidente della Repubblica);

\item {} 
\sphinxAtStartPar
\sphinxstylestrong{5} dal Parlamento in seduta comune;

\item {} 
\sphinxAtStartPar
\sphinxstylestrong{5} dalle supreme magistrature (3 dalla Cassazione, 1 dal Consiglio di Stato, 1 dalla Corte dei Conti).

\end{itemize}

\sphinxAtStartPar
\sphinxstylestrong{Durata e carica:}

\sphinxAtStartPar
I giudici restano in carica \sphinxstylestrong{9 anni} e non possono essere rieletti.

\sphinxAtStartPar
Al proprio interno eleggono un Presidente che rappresenta la Corte.

\sphinxAtStartPar
\sphinxstylestrong{Competenze:}

\sphinxAtStartPar
Ha competenza su 3 materie principali:
\begin{itemize}
\item {} 
\sphinxAtStartPar
È il giudice delle leggi: controlla la \sphinxstylestrong{legittimità costituzionale delle leggi}.

\item {} 
\sphinxAtStartPar
Decide sui \sphinxstylestrong{conflitti di attribuzione tra gli organi dello Stato}.

\item {} 
\sphinxAtStartPar
Decide sull”\sphinxstylestrong{ammissibilità del referendum abrogativo}.

\end{itemize}

\sphinxAtStartPar
Le decisioni sono prese a \sphinxstylestrong{maggioranza.}

\sphinxAtStartPar
\sphinxstylestrong{Il controllo di legittimità costituzionale}

\sphinxAtStartPar
Il compito della Corte è garantire la \sphinxstylestrong{«leggibilità» della Costituzione.}
Il Parlamento ha il potere legislativo, ma se approva una legge in contrasto con la Costituzione, e nemmeno il PdR in sede di rinvio interviene, la Corte può accertare il contrasto.

\sphinxAtStartPar
Questo accertamento avviene attraverso due vie:
\begin{enumerate}
\sphinxsetlistlabels{\arabic}{enumi}{enumii}{}{.}%
\item {} 
\sphinxAtStartPar
\sphinxstylestrong{Via principale (o diretta)}: quando lo Stato o una Regione ritengono che una legge (statale o regionale) invada la propria sfera di competenza e decidono di sottoporre la questione alla Corte.

\item {} 
\sphinxAtStartPar
\sphinxstylestrong{Via incidentale (o indiretta)}: quando, durante un processo di qualunque tipo, un giudice (chiamato «giudice a quo») deve applicare una legge che gli sembra incostituzionale. In questo caso il giudice sospende il giudizio, non applica la legge e rimette la questione alla Corte Costituzionale.

\end{enumerate}

\sphinxstepscope


\chapter{Le altre competenze della Corte Costituzionale}
\label{\detokenize{Le altre competenze della Corte Costituzionale:le-altre-competenze-della-corte-costituzionale}}\label{\detokenize{Le altre competenze della Corte Costituzionale::doc}}
\sphinxAtStartPar
Oltre al controllo di legittimità sulle leggi, la Corte interviene in altri casi fondamentali per l’equilibrio
dello Stato:

\sphinxAtStartPar
\sphinxstylestrong{Conflitti di attribuzione}: Può accadere che all’interno dello Stato si verifichino dei conflitti in cui due
organi rivendicano lo stesso potere.
In pratica, se c’è un conflitto tra organi, ci si può rivolgere alla Corte per chiedere chi ha la competenza.
La Corte decide con una sentenza che è vincolante per le parti.

\sphinxAtStartPar
\sphinxstylestrong{Ammissibilità del referendum}: Il referendum abrogativo non può essere richiesto per tutte le leggi
(sono escluse quelle fiscali, quelle internazionali e di bilancio).
Per evitare che l’abrogazione provochi il blocco dell’attività dello Stato, la richiesta di referendum deve passare
all’esame della Corte, che la deve dichiarare ammissibile.
Se lo fa, il referendum si può svolgere, altrimenti no.

\sphinxAtStartPar
\sphinxstylestrong{Giudizio sul PDR}: Oltre ai tre poteri fondamentali, la Corte può giudicare il PdR (Presidente della Repubblica)
per i reati di alto tradimento o attentato alla Costituzione.
Se lo ritiene colpevole, lo condanna (fino all’ergastolo) e lo rimuove dalle funzioni.
In questo caso la Corte viene integrata: ai 15 giudici ordinari si aggiungono 16 «giudici popolari» (cittadini maggiorenni estratti a sorte tra un elenco compilato dal Parlamento).

\sphinxstepscope


\chapter{Esito del giudizio della Corte}
\label{\detokenize{Esito del giudizio della Corte:esito-del-giudizio-della-corte}}\label{\detokenize{Esito del giudizio della Corte::doc}}
\sphinxAtStartPar
Quando la questione giunge alla Corte, essa ha due possibilità di decisione:

\sphinxAtStartPar
\sphinxstylestrong{Sentenza di rigetto}: La Corte ritiene che la legge non sia in contrasto con la Costituzione.
In questo caso, il giudice del processo originale deve applicare quella legge per risolvere la causa.

\sphinxAtStartPar
\sphinxstylestrong{Sentenza di accoglimento}: La Corte dichiara che la legge è effettivamente incostituzionale.
La sentenza viene pubblicata sulla Gazzetta Ufficiale.
Da quel momento, la legge cessa di avere efficacia e viene eliminata dall’ordinamento giuridico (non può più essere applicata da nessuno).
Sintesi delle competenze (Parte finale)

\sphinxAtStartPar
\sphinxstylestrong{Conflitto di attribuzione}: Se sorge un conflitto tra organi dello Stato su chi debba esercitare un potere,
la Corte decide a chi spetta la competenza.

\sphinxAtStartPar
\sphinxstylestrong{Referendum}: Il referendum abrogativo non è sempre ammesso (es. per leggi tributarie o di bilancio).
La Corte ne valuta l’ammissibilità: se è ammissibile, si procede al voto, altrimenti la richiesta decade.

\sphinxAtStartPar
\sphinxstylestrong{Giudizio sul PdR}: La Corte giudica il Presidente della Repubblica in caso di alto tradimento o attentato alla Costituzione. In questa funzione specifica, la composizione della Corte viene integrata da 16 giudici popolari estratti a sorte.

\sphinxstepscope


\chapter{Il Presidente della Repubblica (PdR)}
\label{\detokenize{Il Presidente della Repubblica:il-presidente-della-repubblica-pdr}}\label{\detokenize{Il Presidente della Repubblica::doc}}
\sphinxAtStartPar
In Italia il PdR è soprattutto un organo di garanzia che tutela la Costituzione, rappresenta l’unità nazionale e
non può prendere decisioni sulla vita politica, che spetta al Parlamento.
Viene eletto dal Parlamento in seduta comune con i rappresentanti delle regioni.
\begin{itemize}
\item {} 
\sphinxAtStartPar
\sphinxstylestrong{Carica}: dura 7 anni e può essere rieletto.

\item {} 
\sphinxAtStartPar
\sphinxstylestrong{Sede}: ha sede a Roma al Quirinale.

\item {} 
\sphinxAtStartPar
\sphinxstylestrong{Termine del mandato}: quando cessa l’incarico diventa Senatore a vita.

\end{itemize}

\sphinxAtStartPar
\sphinxstylestrong{I Poteri del PdR}

\sphinxAtStartPar
Sebbene non abbia poteri politici diretti, la Costituzione gli attribuisce poteri suddivisi in:
\begin{enumerate}
\sphinxsetlistlabels{\arabic}{enumi}{enumii}{}{.}%
\item {} 
\sphinxAtStartPar
\sphinxstylestrong{Poteri Formali}: Rappresenta l’Italia all’estero e nei rapporti internazionali.

\item {} 
\sphinxAtStartPar
\sphinxstylestrong{Poteri Sostanziali}: Poteri effettivi che gli appartengono, tra cui:

\item {} 
\sphinxAtStartPar
\sphinxstylestrong{Promulgare la legge}: prima della scadenza (se la legge è approvata dal Parlamento).

\item {} 
\sphinxAtStartPar
\sphinxstylestrong{Mettere il veto}: può rimandare la legge alle Camere con un messaggio motivato se ritiene che non debba essere approvata (potere di rinvio).

\item {} 
\sphinxAtStartPar
\sphinxstylestrong{Sciogliere il Parlamento}: può decidere di sciogliere le Camere prima della scadenza naturale della legislatura.

\item {} 
\sphinxAtStartPar
\sphinxstylestrong{Nomina}: nomina il Presidente del Consiglio (PdC) e i ministri.

\item {} 
\sphinxAtStartPar
\sphinxstylestrong{Presidenza}: presiede il Consiglio Superiore della Magistratura (CSM).

\item {} 
\sphinxAtStartPar
\sphinxstylestrong{Comando}: è il Capo delle Forze Armate.

\item {} 
\sphinxAtStartPar
\sphinxstylestrong{Nomine speciali}: nomina un terzo (5) dei membri della Corte Costituzionale e nomina 5 Senatori a vita.

\end{enumerate}

\sphinxstepscope


\chapter{Intermediari finanziari}
\label{\detokenize{Intermediari finanziari:intermediari-finanziari}}\label{\detokenize{Intermediari finanziari::doc}}
\sphinxAtStartPar
I risparmiatori hanno il problema dell’inflazione (capitale che perde il proprio potere d’acquisto).
Per evitare questo, l’unico strumento è quello di investire i capitali, operazione che avviene tramite gli intermediari finanziari che operano nel mercato (di capitali).

\sphinxAtStartPar
Per poter svolgere questa attività devono avere dei requisiti:
\begin{enumerate}
\sphinxsetlistlabels{\arabic}{enumi}{enumii}{}{.}%
\item {} 
\sphinxAtStartPar
Essere una S.p.A.: per offrire maggiori garanzie ai risparmiatori.

\item {} 
\sphinxAtStartPar
Essere iscritti in un apposito elenco.

\item {} 
\sphinxAtStartPar
Essere sottoposti al controllo da parte della CONSOB (organismo che controlla la borsa e il mercato finanziario) e della Banca d’Italia.

\end{enumerate}

\sphinxAtStartPar
\sphinxstylestrong{Il Mercato e la Borsa}

\sphinxAtStartPar
La Borsa è il luogo dove avvengono le trattative; qui si scambiano i titoli del mondo e si stabilisce il valore delle monete, che dipende dalla loro domanda e offerta. Alla fine di ogni giorno, la Borsa pubblica i valori per ogni mercato e ogni moneta.

\sphinxAtStartPar
I titoli che si possono acquistare o vendere in borsa sono:
1. \sphinxstylestrong{Le Azioni}: rappresentano una quota della società (S.p.A.); il loro valore dipende dall’andamento della società stessa.
2. \sphinxstylestrong{I Titoli di Stato}: il loro valore dipende dall’andamento economico nazionale.

\sphinxstepscope


\chapter{Mercato aperto dei capitali}
\label{\detokenize{Mercato aperto dei capitali:mercato-aperto-dei-capitali}}\label{\detokenize{Mercato aperto dei capitali::doc}}
\sphinxAtStartPar
Come qualunque altro bene, anche i capitali si vendono:
il mercato dei capitali è il luogo dove si scambiano denaro e titoli.

\sphinxAtStartPar
Questo luogo si chiama Borsa Valori, dove avvengono le contrattazioni tra chi domanda e chi offre capitale.

\sphinxAtStartPar
In Italia la borsa principale ha sede a Milano.
In questo luogo, alla fine di ogni seduta, viene stabilito il valore di ogni moneta e titolo, permettendo di confrontarne l’andamento giornaliero (che può essere al rialzo o al ribasso).

\sphinxAtStartPar
Il mercato monetario si divide in:
\begin{enumerate}
\sphinxsetlistlabels{\arabic}{enumi}{enumii}{}{.}%
\item {} 
\sphinxAtStartPar
\sphinxstylestrong{Mercato primario}: quello in cui si vendono i titoli di nuova emissione; è utilizzato dalle società o dallo Stato per raccogliere capitali.

\item {} 
\sphinxAtStartPar
\sphinxstylestrong{Mercato secondario}: quello in cui vengono negoziati i titoli già in circolazione.

\end{enumerate}

\sphinxAtStartPar
\sphinxstylestrong{Tipologie di Titoli}

\sphinxAtStartPar
I titoli si dividono principalmente in:
\begin{enumerate}
\sphinxsetlistlabels{\arabic}{enumi}{enumii}{}{.}%
\item {} 
\sphinxAtStartPar
\sphinxstylestrong{Titoli Pubblici} (BOT, CCT): emessi dallo Stato. Hanno la caratteristica di avere tassi di interesse generalmente bassi, ma sono considerati investimenti sicuri poiché lo Stato ne garantisce il rimborso.

\item {} 
\sphinxAtStartPar
\sphinxstylestrong{Titoli Privati} (Azioni e Obbligazioni): emessi dalle S.p.A. Le obbligazioni consistono in prestiti alla società con un tasso di interesse variabile.

\end{enumerate}

\sphinxAtStartPar
\sphinxstylestrong{Elementi del mercato finanziario}
\begin{itemize}
\item {} 
\sphinxAtStartPar
\sphinxstylestrong{Strumenti finanziari}: comprendono la moneta estera, i titoli e i «futures» (contratti d’investimento).

\item {} 
\sphinxAtStartPar
\sphinxstylestrong{Operatori di Borsa}: sono professionisti abilitati (come i broker) che possono operare all’interno della borsa per conto dei clienti.

\item {} 
\sphinxAtStartPar
\sphinxstylestrong{Andamento del mercato}: dipende dalla domanda e dall’offerta. Se la domanda di una moneta o di un titolo aumenta, il suo valore sale; se l’offerta è eccessiva, il valore scende. Anche le decisioni della Banca Centrale sugli interessi influenzano l’andamento dei mercati.

\end{itemize}

\sphinxstepscope


\chapter{Il Sistema Bancario}
\label{\detokenize{Il Sistema Bancario:il-sistema-bancario}}\label{\detokenize{Il Sistema Bancario::doc}}
\sphinxAtStartPar
Le banche costituiscono l’elemento essenziale del sistema economico.
Infatti, esse rappresentano il collegamento tra coloro che offrono risparmio (le famiglie) e coloro che chiedono capitale (le imprese). Per questa ragione, le banche, nel momento in cui offrono prestiti, si fanno pagare un tasso d’interesse che deve essere superiore a quello che hanno pagato per avere il denaro.

\sphinxAtStartPar
\sphinxstylestrong{La Banca Centrale e il costo del denaro}

\sphinxAtStartPar
Normalmente le banche prestano il denaro dei risparmiatori e quello che ottengono dalla Banca Centrale. Ogni Stato ha una sua Banca Centrale (in Italia è la Banca d’Italia, che a sua volta dipende dalla BCE \sphinxhyphen{} Banca Centrale Europea), la quale stabilisce il costo del denaro (Tasso Ufficiale).

\sphinxAtStartPar
\sphinxstyleemphasis{L’andamento dell’economia è influenzato da queste variazioni:}
\begin{itemize}
\item {} 
\sphinxAtStartPar
\sphinxstylestrong{Aumento dei tassi}: Se il tasso ufficiale e quindi il costo del denaro cresce, anche le banche dovranno applicare ai propri clienti un tasso superiore. In questo modo si scoraggiano gli investimenti.

\item {} 
\sphinxAtStartPar
\sphinxstylestrong{Diminuzione dei tassi}: Al contrario, se si diminuisce il tasso pagato, questo provocherà un aumento degli investimenti e dei consumi.

\end{itemize}

\sphinxAtStartPar
\sphinxstylestrong{Politica monetaria e Inflazione}

\sphinxAtStartPar
La Banca Centrale agisce sull’andamento dell’economia in base al rischio di inflazione:
\begin{enumerate}
\sphinxsetlistlabels{\arabic}{enumi}{enumii}{}{.}%
\item {} 
\sphinxAtStartPar
Se non ci sono pericoli d’inflazione, può abbassare il tasso ufficiale per stimolare l’economia.

\item {} 
\sphinxAtStartPar
Al contrario, se c’è pericolo di inflazione, deve alzare il tasso ufficiale per frenare la circolazione monetaria.

\end{enumerate}

\sphinxstepscope


\chapter{Gli operatori del mercato finanziario}
\label{\detokenize{Gli operatori del mercato finanziario:gli-operatori-del-mercato-finanziario}}\label{\detokenize{Gli operatori del mercato finanziario::doc}}
\sphinxAtStartPar
In base al mandato con cui il cliente richiede di acquistare o vendere strumenti finanziari per effettuare l’operazione, i soggetti che operano sul mercato possono essere divisi in due grandi categorie:
\begin{enumerate}
\sphinxsetlistlabels{\arabic}{enumi}{enumii}{}{.}%
\item {} 
\sphinxAtStartPar
\sphinxstylestrong{I Risparmiatori (o Cassettisti)}: rappresentano i soggetti prudenti. Il loro obiettivo è proteggere il proprio capitale nel tempo, senza correre rischi eccessivi, cercando un rendimento costante.

\item {} \begin{description}
\sphinxlineitem{\sphinxstylestrong{Gli Speculatori}: sono soggetti che vogliono rischiare per ottenere guadagni elevati in tempi brevi, sfruttando le variazioni di prezzo dei titoli.}\begin{itemize}
\item {} \begin{description}
\sphinxlineitem{Gli speculatori, a loro volta, si dividono in base alla strategia adottata e alle previsioni sull’andamento del mercato:}\begin{enumerate}
\sphinxsetlistlabels{\arabic}{enumii}{enumiii}{}{.}%
\item {} 
\sphinxAtStartPar
\sphinxstylestrong{Speculatori Rialzisti (Toro)}: acquistano titoli oggi scommettendo che il loro valore aumenterà in futuro, per poi rivenderli e guadagnare sulla differenza di prezzo.

\item {} 
\sphinxAtStartPar
\sphinxstylestrong{Speculatori Ribassisti (Orso)}: vendono titoli (spesso non ancora posseduti fisicamente) scommettendo che il loro valore diminuirà, per poi riacquistarli a un prezzo inferiore e guadagnare sul calo.

\end{enumerate}

\end{description}

\end{itemize}

\end{description}

\end{enumerate}

\sphinxAtStartPar
\sphinxstylestrong{Il valore dei titoli e la domanda}

\sphinxAtStartPar
Il valore di un titolo o di una moneta è strettamente legato alla quantità domandata e offerta sul mercato.

\sphinxAtStartPar
Ragione per cui, se lo Stato o un’azienda è in una situazione economica solida, la domanda per i suoi titoli aumenta e,
di conseguenza, ne cresce il valore rispetto ad altre opzioni di investimento.

\sphinxstepscope


\chapter{Il Mercato dei cambi}
\label{\detokenize{Il Mercato dei cambi:il-mercato-dei-cambi}}\label{\detokenize{Il Mercato dei cambi::doc}}
\sphinxAtStartPar
Ogni moneta ha un suo valore, che rappresenta la quantità di beni e servizi che si possono acquistare con una
determinata pressione monetaria.

\sphinxAtStartPar
Nel momento in cui si vuole acquistare la moneta di un altro Stato, o si vogliono fare acquisti in uno Stato diverso che ha una moneta diversa, si parla di tasso di cambio.

\sphinxAtStartPar
\sphinxstylestrong{Il tasso di cambio} è il prezzo a cui una moneta viene scambiata con un’altra.

\sphinxAtStartPar
\sphinxstylestrong{Determinazione del valore}

\sphinxAtStartPar
Il valore di ogni moneta varia quotidianamente ed è determinato dal mercato in base alla domanda e all’offerta.
\begin{enumerate}
\sphinxsetlistlabels{\arabic}{enumi}{enumii}{}{.}%
\item {} 
\sphinxAtStartPar
Se molti soggetti richiedono una determinata valuta (per investimenti o commercio), il suo \sphinxstylestrong{valore sale}.

\item {} 
\sphinxAtStartPar
Se c’è eccesso di offerta, il \sphinxstylestrong{valore scende}.

\end{enumerate}

\sphinxAtStartPar
Nella Borsa, oltre ai titoli, vengono quotidianamente scambiate e quotate anche le monete dei vari Stati,
permettendo di stabilire il rapporto di valore tra di esse (es. Euro/Dollaro).



\renewcommand{\indexname}{Indice}
\printindex
\end{document}